\documentclass{article}
\usepackage{amsmath,amssymb,amsthm}
\usepackage{booktabs}
\usepackage{hyperref}
\usepackage{listings}
\usepackage{xcolor}
\usepackage[margin=1in]{geometry}
\usepackage{array}
\usepackage{multirow}

\newcommand{\hex}[1]{\texttt{\$#1}}
\newcommand{\opcode}[1]{\texttt{#1}}

\lstset{
  basicstyle=\ttfamily\small,
  columns=fullflexible,
  frame=single,
  backgroundcolor=\color{gray!10},
}

\title{Deep Learning Approaches for 6502 Self-Replicator\\Generation and Analysis: A Survey}
\author{}
\date{March 2026}

\begin{document}
\maketitle

\begin{abstract}
We survey deep learning approaches applicable to the generation, classification,
and optimization of self-replicating programs in a 6502-based cellular automaton.
Our setting involves short byte sequences (8--32 bytes) over a 256-token vocabulary,
where viable replicators occur with probability $\approx 2^{-61}$ under uniform
sampling.  An existing WFST (weighted finite-state transducer) pipeline provides
exact structural inference but is limited to finite-state representations.
We examine seven complementary neural approaches---autoregressive generation,
surrogate classification, reinforcement learning, sequence-to-sequence mutation,
masked prediction, state-space models, and hybrid WFST+neural architectures---analyzing
their architecture, data requirements, and trade-offs against the WFST baseline.
\end{abstract}

\tableofcontents

%% ============================================================
\section{Problem Setting}
\label{sec:setting}

The 6502life world consists of a $256 \times 256$ grid of cells, each running a
6502 CPU with 1024 bytes of memory.  A \emph{self-replicating program} is a byte
sequence that, when loaded into a cell and executed under the preemptive scheduler,
copies itself into at least one neighboring cell.

\paragraph{Minimal replicator structure.}
The canonical minimal replicator is an 8-byte loop:
\begin{lstlisting}
  B5 00   LDA $00,X    ; load from zero-page + X
  9D 00 04  STA $0400,X  ; store to neighbor base + X
  E8      INX          ; increment index
  90 F8   BCC loop     ; branch if X < 256
\end{lstlisting}
followed by a \opcode{BRK} with appropriate operand to trigger the copy/swap
syscall.  Functional equivalence classes allow variations in addressing mode,
branch type, target cell, and surrounding NOP slides.

\paragraph{WFST baseline.}
Our existing pipeline models byte sequences as paths through a weighted
finite-state transducer whose transition weights are trained from simulation
oracle labels via importance sampling.  The WFST computes
\begin{equation}
  P(\text{viable} \mid L) = \sum_{\mathbf{x} \in \{0,\ldots,255\}^L}
    w_{\text{WFST}}(\mathbf{x}) \cdot \mathbf{1}[\text{oracle}(\mathbf{x})=1],
  \label{eq:wfst-prob}
\end{equation}
yielding $P(\text{viable} \mid L) \approx 2^{-61}$ for the 1024-byte cell.
The WFST provides exact inference over its state space but cannot capture
dependencies beyond its finite memory (e.g., interactions between a \opcode{SEC}
at position $i$ and a \opcode{BCS} at position $j \gg i$).

\paragraph{Notation.}
We denote a byte sequence of length $L$ as
$\mathbf{x} = (x_1, \ldots, x_L) \in \mathcal{V}^L$ where $|\mathcal{V}|=256$.
The simulation oracle is $f:\mathcal{V}^L \to \{0,1\}$ (binary viability)
or $f:\mathcal{V}^L \to \mathbb{R}_{\geq 0}$ (spread count).
A single oracle call requires ${\sim}10^4$ scheduler cycles (${\sim}1$\,ms on CPU).


%% ============================================================
\section{Autoregressive Language Model}
\label{sec:lm}

\subsection{Architecture}

An autoregressive transformer models the joint distribution over byte sequences:
\begin{equation}
  p_\theta(\mathbf{x}) = \prod_{t=1}^{L} p_\theta(x_t \mid x_{<t}).
  \label{eq:ar}
\end{equation}
Each token $x_t \in \{0,\ldots,255\}$ is embedded into $d$-dimensional space.
Given the short sequence lengths ($L \leq 32$), a small model suffices:
\begin{itemize}
  \item Embedding dimension $d = 64$--$128$
  \item 2--4 transformer layers, 2--4 attention heads
  \item Causal (left-to-right) attention mask
  \item Vocabulary: 256 byte values + start/end tokens
  \item Positional encoding: learned or sinusoidal (sequence is short enough
        that both work)
  \item Total parameters: ${\sim}50$K--$500$K
\end{itemize}

At inference, we sample $\mathbf{x} \sim p_\theta$ autoregressively with
temperature control, top-$k$, or nucleus sampling.

\subsection{Training}

\paragraph{Data.}
Training requires a corpus of viable replicator sequences.  Since viable
replicators are extremely rare under uniform sampling ($2^{-61}$), we use
the WFST importance sampler to generate candidates, then filter through the
simulation oracle.  With WFST proposal efficiency of ${\sim}2^{-10}$
relative to uniform, we need ${\sim}2^{10} \approx 1000$ WFST samples
per viable sequence, so ${\sim}10^6$ WFST draws yield ${\sim}10^3$ viable
training examples.

\paragraph{Objective.}
Standard cross-entropy loss on viable sequences:
\begin{equation}
  \mathcal{L}(\theta) = -\mathbb{E}_{\mathbf{x} \sim \mathcal{D}_{\text{viable}}}
    \left[\sum_{t=1}^{L} \log p_\theta(x_t \mid x_{<t})\right].
\end{equation}
Optionally, weight training examples by spread count to bias toward
high-fitness replicators.

\subsection{Advantages over WFST}

\begin{itemize}
  \item \textbf{Long-range dependencies.}  Self-attention captures interactions
        between arbitrary positions (e.g., a \opcode{SEC} at byte 2 enabling
        \opcode{BCS} at byte 9), while the WFST is limited to its state memory.
  \item \textbf{Soft generalization.}  The transformer can interpolate between
        training examples, potentially generating novel replicator families not
        in the WFST's equivalence classes.
  \item \textbf{Scalability.}  Adding new replicator families requires only
        more training data, not manual transducer engineering.
\end{itemize}

\subsection{Limitations}

\begin{itemize}
  \item No guarantee of structural validity---generated sequences may be
        syntactically plausible but non-functional.
  \item Requires a substantial corpus of viable examples (thousands), which
        are expensive to obtain.
  \item Cannot provide exact probabilities or guarantees, unlike the WFST.
  \item For 8-byte sequences, the WFST may already capture the relevant
        structure; the transformer's advantage grows with sequence length.
\end{itemize}

\subsection{Related Work}

Autoregressive models over byte-level program representations have been
explored in neural program synthesis~\cite{chen2021codex}, though typically
for high-level languages.  The closest analogue is byte-level language
modeling for binary analysis~\cite{pei2020trex}, which applies transformers
to disassembled binary code for function similarity detection.


%% ============================================================
\section{Surrogate Classifier (Learned Oracle)}
\label{sec:classifier}

\subsection{Architecture}

A discriminative model approximates the simulation oracle:
\begin{equation}
  \hat{f}_\phi(\mathbf{x}) = \sigma\bigl(g_\phi(\mathbf{x})\bigr)
  \approx P(\text{viable} \mid \mathbf{x}),
\end{equation}
where $g_\phi$ is a neural network and $\sigma$ is the sigmoid function.

\paragraph{Option A: 1D CNN.}
Treat the byte sequence as a 1D signal with 256-dimensional one-hot encoding
(or learned byte embeddings):
\begin{itemize}
  \item Embedding: $256 \to 64$
  \item Conv layers: 3 layers with kernels of size 3, 5, 7 (capturing
        opcode-operand structure), 64--128 filters each
  \item Global average pooling $\to$ FC $\to$ sigmoid
  \item Parameters: ${\sim}100$K
\end{itemize}
The CNN naturally captures local opcode-operand groupings (1--3 byte
instructions) through its receptive field.

\paragraph{Option B: Small transformer encoder.}
\begin{itemize}
  \item Byte embedding $d=64$, 2 encoder layers, 2 heads
  \item \texttt{[CLS]} token pooling $\to$ FC $\to$ sigmoid
  \item Parameters: ${\sim}50$K
\end{itemize}

\subsection{Training}

\paragraph{Data.}
Binary classification requires both positive (viable) and negative
(non-viable) examples.  The extreme class imbalance ($2^{-61}$) means
uniform negatives are trivially classifiable.  Instead, we construct
\emph{hard negatives} via:
\begin{enumerate}
  \item WFST samples that pass structural checks but fail simulation
        (``near-miss'' programs)
  \item Single-byte mutations of viable replicators
  \item Random programs seeded with correct opcodes but wrong operands
\end{enumerate}
A balanced dataset of ${\sim}10^4$ positive and ${\sim}10^4$ hard negative
examples suffices for the initial classifier.

\paragraph{Objective.}
Binary cross-entropy with optional focal loss~\cite{lin2017focal} to handle
residual class imbalance:
\begin{equation}
  \mathcal{L}(\phi) = -\mathbb{E}\bigl[
    \alpha_t (1 - \hat{f}_\phi)^\gamma \log \hat{f}_\phi
    + (1 - \alpha_t) \hat{f}_\phi^\gamma \log(1 - \hat{f}_\phi)
  \bigr].
\end{equation}

\subsection{Application}

The surrogate classifier serves as a \emph{fast filter} in the sampling
pipeline:
\begin{equation}
  \text{WFST sample} \xrightarrow{\hat{f}_\phi > \tau}
  \text{simulation oracle} \xrightarrow{f=1} \text{accept}.
\end{equation}
If $\hat{f}_\phi$ achieves 95\% recall at 10$\times$ precision improvement,
simulation cost drops by $10\times$.  The classifier evaluates in
${\sim}10\,\mu$s (GPU-batched) vs.\ ${\sim}1$\,ms for simulation.

\subsection{Advantages over WFST}

\begin{itemize}
  \item Orders of magnitude faster than simulation for filtering.
  \item Can learn nonlinear viability criteria that the WFST's finite
        states cannot represent.
  \item Naturally handles soft scoring (confidence), enabling ranking.
\end{itemize}

\subsection{Limitations}

\begin{itemize}
  \item False negatives discard viable replicators; false positives waste
        simulation budget.
  \item The decision boundary may not generalize to replicator families
        absent from training data.
  \item Requires periodic retraining as new replicator types are discovered.
\end{itemize}

\subsection{Related Work}

Surrogate models for expensive simulations are standard in Bayesian
optimization~\cite{snoek2012practical} and materials science.
Neural network surrogates for program behavior prediction appear in
learned interpreters~\cite{bieber2020learning} and neural execution
engines.


%% ============================================================
\section{Reinforcement Learning (PPO/REINFORCE)}
\label{sec:rl}

\subsection{Formulation}

We cast replicator generation as a sequential decision problem:
\begin{itemize}
  \item \textbf{State} $s_t = (x_1, \ldots, x_{t-1})$: bytes generated so far
  \item \textbf{Action} $a_t = x_t \in \{0,\ldots,255\}$: next byte
  \item \textbf{Policy} $\pi_\theta(a_t \mid s_t)$: autoregressive model
        (Section~\ref{sec:lm})
  \item \textbf{Reward} $R(\mathbf{x})$: simulation spread count (or binary
        viability), received at episode end ($t=L$)
\end{itemize}

\paragraph{REINFORCE.}
The policy gradient estimator:
\begin{equation}
  \nabla_\theta J = \mathbb{E}_{\mathbf{x} \sim \pi_\theta}
    \left[R(\mathbf{x}) \sum_{t=1}^{L} \nabla_\theta \log \pi_\theta(x_t \mid x_{<t})\right].
  \label{eq:reinforce}
\end{equation}
This is structurally identical to the WFST training loop: both use
likelihood-ratio gradient estimation with oracle reward.  The difference
is the policy class ($\pi_\theta$ is a neural network vs.\ a WFST).

\paragraph{PPO.}
Proximal Policy Optimization~\cite{schulman2017ppo} stabilizes training via
a clipped surrogate objective:
\begin{equation}
  \mathcal{L}^{\text{PPO}}(\theta) = \mathbb{E}\left[
    \min\bigl(r_t(\theta) \hat{A}_t,\;
    \text{clip}(r_t(\theta), 1{-}\epsilon, 1{+}\epsilon)\hat{A}_t\bigr)
  \right],
\end{equation}
where $r_t(\theta) = \pi_\theta(a_t \mid s_t) / \pi_{\theta_{\text{old}}}(a_t \mid s_t)$
and $\hat{A}_t$ is the advantage estimate.  For terminal-reward-only MDPs,
the advantage is constant across timesteps within an episode, simplifying to
$\hat{A}_t = R(\mathbf{x}) - b$ where $b$ is a baseline (e.g., running mean reward).

\subsection{Architecture}

The policy network is the autoregressive model from Section~\ref{sec:lm}.
PPO additionally requires a value network $V_\psi(s_t)$; for terminal
rewards, a single scalar baseline suffices.

\subsection{Training}

Each training iteration:
\begin{enumerate}
  \item Sample $N$ byte sequences from $\pi_\theta$
  \item Evaluate each via simulation oracle ($N$ simulation runs)
  \item Compute policy gradient and update $\theta$
\end{enumerate}

\paragraph{Sample efficiency.}
With reward signal $\approx 2^{-61}$ under uniform policy, pure RL from
scratch is intractable.  Two mitigations:
\begin{enumerate}
  \item \textbf{Warm-start}: pretrain $\pi_\theta$ on viable examples
        (Section~\ref{sec:lm}), then fine-tune with RL. The pretrained
        policy concentrates probability mass near viable sequences.
  \item \textbf{Reward shaping}: use the WFST structural score or the
        surrogate classifier (Section~\ref{sec:classifier}) as a dense
        reward signal, supplementing the sparse simulation reward.
\end{enumerate}
With warm-start, a viable sequence appears roughly every ${\sim}10^3$
samples, making PPO feasible with ${\sim}10^5$ total simulation calls.

\subsection{Advantages over WFST}

\begin{itemize}
  \item \textbf{Richer policy class}: neural policies can represent
        multimodal distributions over replicator families, while the
        WFST is limited to its state topology.
  \item \textbf{Continuous optimization}: PPO/REINFORCE optimize in
        continuous parameter space via gradient descent, vs.\ the WFST's
        discrete structure.
  \item \textbf{Multi-objective}: reward can incorporate spread speed,
        mutation robustness, compactness, etc.\ without redesigning the
        transducer.
\end{itemize}

\subsection{Limitations}

\begin{itemize}
  \item High variance in gradient estimates for rare-event rewards.
  \item Simulation cost: each policy evaluation requires a full board
        simulation.
  \item Risk of mode collapse to a single replicator family.
  \item The WFST's REINFORCE-like training is simpler and may suffice
        for the current problem scale.
\end{itemize}

\subsection{Related Work}

RL for program synthesis has been explored in
neural program induction~\cite{bunel2018leveraging} and
code generation~\cite{le2022coderl}.
The formulation closely resembles REINFORCE-based text
generation~\cite{williams1992simple,ranzato2016sequence}.


%% ============================================================
\section{Seq2Seq Evolutionary Model}
\label{sec:seq2seq}

\subsection{Architecture}

An encoder-decoder transformer maps a parent replicator to a viable
offspring:
\begin{equation}
  p_\theta(\mathbf{x}' \mid \mathbf{x}) = \prod_{t=1}^{L}
    p_\theta(x'_t \mid x'_{<t}, \mathbf{x}),
\end{equation}
where $\mathbf{x}$ is a viable parent and $\mathbf{x}'$ is the mutated
offspring.

\begin{itemize}
  \item Encoder: 2-layer transformer, $d=64$, encodes parent sequence
  \item Decoder: 2-layer transformer with cross-attention to encoder
  \item Shared byte embedding across encoder and decoder
  \item Parameters: ${\sim}100$K--$200$K
\end{itemize}

\subsection{Training Data}

Pairs $(\mathbf{x}, \mathbf{x}')$ where both are viable and $\mathbf{x}'$
is derived from $\mathbf{x}$ via:
\begin{enumerate}
  \item \textbf{Simulation lineage}: run the board, extract parent-child
        pairs from copy events (includes mutation noise from \texttt{pBitNoise}).
  \item \textbf{Synthetic mutations}: apply $k$-bit flips to viable
        sequences, keep pairs where both remain viable.
  \item \textbf{WFST neighborhood}: sample sequences near a viable program
        in WFST weight space.
\end{enumerate}
Target: ${\sim}10^4$ viable pairs.

\subsection{Application}

Given a known replicator $\mathbf{x}$, the model generates a distribution
over viable neighbors.  This enables:
\begin{itemize}
  \item \textbf{Directed evolution}: hill-climb in viability/fitness space
        by iterating the model.
  \item \textbf{Diversity search}: sample multiple offspring per parent to
        explore the replicator landscape.
  \item \textbf{Robustness analysis}: characterize the viable neighborhood
        of a given program.
\end{itemize}

\subsection{Advantages over WFST}

\begin{itemize}
  \item Models the \emph{mutation kernel} directly, rather than the marginal
        distribution over viable sequences.
  \item Can learn which positions tolerate mutation and which are critical.
  \item Preserves functional structure while varying non-essential bytes.
\end{itemize}

\subsection{Limitations}

\begin{itemize}
  \item Requires pairs of viable programs, which are harder to obtain than
        individual viable programs.
  \item May learn to copy the input with minimal changes (identity shortcut).
        Mitigate with minimum edit distance constraints or dropout on the
        encoder.
  \item Does not discover \emph{de novo} replicator families---only
        transforms known ones.
\end{itemize}

\subsection{Related Work}

Program repair and transformation using seq2seq models is established
in software engineering~\cite{tufano2019empirical}.  The evolutionary
framing connects to learned mutation operators in
genetic programming~\cite{chen2020program}.


%% ============================================================
\section{Masked Bidirectional Model}
\label{sec:masked}

\subsection{Architecture}

A BERT-style~\cite{devlin2019bert} masked language model trained on viable
replicator sequences.  During training, a random subset of bytes is masked
and the model predicts them from bidirectional context:
\begin{equation}
  \mathcal{L}(\theta) = -\mathbb{E}_{\mathbf{x} \sim \mathcal{D}}
    \left[\sum_{t \in \mathcal{M}} \log p_\theta(x_t \mid \mathbf{x}_{\setminus\mathcal{M}})\right],
\end{equation}
where $\mathcal{M} \subset \{1,\ldots,L\}$ is the set of masked positions
(typically 15\% of tokens).

\begin{itemize}
  \item Encoder-only transformer, 2--4 layers, $d=64$, 2--4 heads
  \item Special \texttt{[MASK]} token in the 257-token vocabulary
  \item Parameters: ${\sim}50$K--$200$K
\end{itemize}

\subsection{Training}

Same viable-sequence corpus as Section~\ref{sec:lm} (${\sim}10^3$--$10^4$
examples).  Data augmentation via:
\begin{itemize}
  \item NOP slide insertion/deletion (shifting positions without changing
        semantics)
  \item Equivalent opcode substitution (e.g., \opcode{INX} $\leftrightarrow$
        \opcode{INY} with matching address mode changes)
\end{itemize}

\subsection{Application}

\paragraph{Credit assignment.}
Given a partial program with unknown bytes, the model fills them in:
``What byte at position 5 makes this a viable replicator, given the
surrounding context?''  This directly addresses the credit assignment
problem---which bytes are critical and which are free.

\paragraph{Iterative refinement.}
Starting from a random sequence, iteratively mask and re-predict positions,
annealing the mask rate.  This Gibbs-like sampling procedure can converge
to viable sequences without autoregressive left-to-right generation.

\paragraph{NOP slide design.}
Mask the slide region and let the model predict which byte values serve
as functional NOPs given the surrounding loop structure.

\subsection{Advantages over WFST}

\begin{itemize}
  \item \textbf{Bidirectional context}: can condition on both preceding
        and following bytes, unlike the WFST's left-to-right processing.
  \item \textbf{Natural infilling}: directly answers ``what goes here?''
        queries without marginalization.
  \item \textbf{Credit assignment}: per-position confidence scores indicate
        which bytes are constrained vs.\ free.
\end{itemize}

\subsection{Limitations}

\begin{itemize}
  \item No tractable joint likelihood $p_\theta(\mathbf{x})$---cannot
        directly estimate replicator probability.
  \item Iterative sampling (Gibbs) is slower than single-pass autoregressive
        generation.
  \item Mask rate and annealing schedule require tuning.
\end{itemize}

\subsection{Related Work}

Masked language models for code have been explored for code
understanding~\cite{feng2020codebert} and completion.
The iterative refinement procedure connects to
non-autoregressive machine translation~\cite{ghazvininejad2019mask}.


%% ============================================================
\section{State-Space Models (Mamba-Style)}
\label{sec:ssm}

\subsection{Motivation}

The WFST is, structurally, a \emph{hand-designed state-space model}:
it maintains a finite hidden state $h_t$ and updates it linearly as it
consumes each byte.  Modern learned state-space models
(SSMs)~\cite{gu2024mamba} generalize this with:
\begin{equation}
  h_t = \bar{A} h_{t-1} + \bar{B} x_t, \qquad
  y_t = C h_t,
\end{equation}
where $\bar{A}, \bar{B}$ are input-dependent (selective) parameters
and $h_t \in \mathbb{R}^N$ is a continuous hidden state.

\subsection{Architecture}

A Mamba-style SSM for byte sequences:
\begin{itemize}
  \item Input: byte embedding, $d=64$
  \item 2--4 Mamba blocks with state dimension $N=16$--$64$
  \item Linear-time complexity: $O(L \cdot d \cdot N)$, matching the WFST
  \item Output head: next-byte prediction (generative) or classification
  \item Parameters: ${\sim}30$K--$150$K
\end{itemize}

\subsection{Relationship to WFST}

The WFST is a special case: its transition matrices are sparse (one nonzero
entry per column in the deterministic case) and its state space is discrete
and finite.  A learned SSM with continuous state $h_t \in \mathbb{R}^N$
can:
\begin{enumerate}
  \item \textbf{Recover} the WFST dynamics if the problem is truly
        finite-state (the SSM can learn to quantize its states).
  \item \textbf{Extend} beyond finite states to capture soft, graded
        structural features---e.g., ``how close is this prefix to needing
        a branch instruction?''
  \item \textbf{Learn nonlinear effects} via the selective (input-dependent)
        gating mechanism, which the WFST's fixed transition weights cannot
        represent.
\end{enumerate}

\subsection{Training}

Same data and objectives as Section~\ref{sec:lm} (autoregressive) or
Section~\ref{sec:classifier} (discriminative).  The SSM replaces the
transformer backbone in either case.

\subsection{Advantages over WFST}

\begin{itemize}
  \item Same linear-time sequential processing, but with learned
        (rather than hand-designed) state dynamics.
  \item Continuous state enables soft generalization beyond the WFST's
        discrete equivalence classes.
  \item Selective gating learns input-dependent transitions, unlike the
        WFST's fixed (input-independent modulo current state) weights.
\end{itemize}

\subsection{Limitations}

\begin{itemize}
  \item For very short sequences ($L \leq 32$), the efficiency advantage
        over attention ($O(L^2)$) is negligible---$32^2 = 1024$ is small.
  \item SSMs may underperform transformers on tasks requiring precise
        positional reasoning (e.g., ``byte 3 must equal byte 7'').
  \item Less interpretable than the WFST, whose states have explicit
        semantic meaning.
\end{itemize}

\subsection{Related Work}

Mamba~\cite{gu2024mamba} and its successors demonstrate strong sequence
modeling with linear complexity.  Application to byte-level code modeling
is unexplored but natural given the sequential structure.


%% ============================================================
\section{Hybrid WFST+Neural Architecture}
\label{sec:hybrid}

\subsection{Motivation}

The WFST excels at hard structural constraints (instruction decoding,
address mode parsing, loop structure detection) that are \emph{exactly
specified} by the 6502 ISA.  Neural networks excel at learning soft
statistical patterns from data.  A hybrid architecture uses each component
for what it does best.

\subsection{Architecture Variants}

\paragraph{Variant A: WFST proposal + neural importance weights.}
The WFST generates candidate sequences $\mathbf{x} \sim q_{\text{WFST}}$.
A neural network computes importance weights:
\begin{equation}
  w(\mathbf{x}) = \frac{p_\phi(\text{viable} \mid \mathbf{x})}
                       {q_{\text{WFST}}(\mathbf{x})},
\end{equation}
where $p_\phi$ is the surrogate classifier (Section~\ref{sec:classifier}).
This importance-weighted sampling corrects the WFST's distributional
errors while retaining its structural guarantees.

\paragraph{Variant B: WFST features + neural classifier.}
Extract the WFST's state trajectory
$\mathbf{h} = (h_1, \ldots, h_L)$ as a feature sequence and feed it
to a neural classifier:
\begin{equation}
  \hat{f}(\mathbf{x}) = \text{NN}(\mathbf{h}_1, \ldots, \mathbf{h}_L).
\end{equation}
The WFST states provide a strong inductive bias (instruction boundaries,
loop detection, etc.) and the neural network learns residual patterns
the WFST misses.

\paragraph{Variant C: Neural-weighted WFST.}
Replace the WFST's scalar transition weights with neural network outputs.
At each position $t$, a small network computes the weight of each
transition given the prefix:
\begin{equation}
  w(h_t \to h_{t+1} \mid x_t, \text{context}_t)
    = \text{MLP}(\text{embed}(h_t), \text{embed}(x_t), c_t),
\end{equation}
where $c_t$ is a neural context vector summarizing $x_{<t}$.
This preserves the WFST's topology (structural constraints as hard
edges) while making weights context-dependent.

\subsection{Training}

All variants use the same simulation-labeled data.  Variant C requires
differentiable WFST forward passes, implementable via the forward
algorithm on the semiring of real-valued weights.

\subsection{Advantages}

\begin{itemize}
  \item \textbf{Best of both worlds}: exact structural constraints from
        the WFST, learned soft patterns from the neural network.
  \item \textbf{Data efficiency}: the WFST provides a strong prior,
        reducing the neural network's learning burden.
  \item \textbf{Interpretability}: the WFST component remains inspectable;
        the neural component handles the ``residual'' complexity.
\end{itemize}

\subsection{Limitations}

\begin{itemize}
  \item Implementation complexity: requires maintaining both WFST and
        neural components.
  \item The WFST topology may constrain the hybrid system if the true
        structure deviates from the hand-designed states.
  \item Gradient flow through the WFST requires careful engineering
        (or approximation in Variant A).
\end{itemize}

\subsection{Related Work}

Hybrid FST+neural models appear in speech recognition, where
WFST-based decoders are combined with neural acoustic
models~\cite{mohri2002weighted}.  Neural-weighted automata have
been studied in formal language theory~\cite{schwartz2018sopa}.


%% ============================================================
\section{Comparison}
\label{sec:comparison}

\begin{table}[htbp]
\centering
\caption{Comparison of approaches for 6502 replicator generation and analysis.
Training data is measured in viable replicator examples (positive class).
Inference speed is per-sequence on GPU unless noted.}
\label{tab:comparison}
\small
\renewcommand{\arraystretch}{1.2}
\begin{tabular}{@{}l c c c c c@{}}
\toprule
\textbf{Approach} &
\textbf{Task} &
\textbf{Params} &
\textbf{Training data} &
\textbf{Inference} &
\textbf{Key advantage} \\
\midrule
WFST (baseline) &
Generate &
${\sim}10^3$ weights &
${\sim}10^3$ oracle calls &
${\sim}1\,\mu$s &
Exact inference \\
\midrule
Autoregressive LM &
Generate &
50K--500K &
${\sim}10^3$ viable &
${\sim}100\,\mu$s &
Long-range deps \\
Surrogate classifier &
Classify &
50K--100K &
${\sim}10^4$ pairs &
${\sim}10\,\mu$s &
100$\times$ faster oracle \\
RL (PPO) &
Generate &
50K--500K &
${\sim}10^5$ sim calls &
${\sim}100\,\mu$s &
Multi-objective \\
Seq2seq &
Mutate &
100K--200K &
${\sim}10^4$ pairs &
${\sim}200\,\mu$s &
Directed evolution \\
Masked (BERT) &
Infill &
50K--200K &
${\sim}10^3$ viable &
${\sim}50\,\mu$s &
Credit assignment \\
SSM (Mamba) &
Gen./Class. &
30K--150K &
${\sim}10^3$ viable &
${\sim}50\,\mu$s &
Learned WFST \\
Hybrid WFST+NN &
Gen./Class. &
50K--200K &
${\sim}10^3$ viable &
${\sim}10\,\mu$s &
Best of both \\
\bottomrule
\end{tabular}
\end{table}

\paragraph{Recommended development order.}
\begin{enumerate}
  \item \textbf{Surrogate classifier} (Section~\ref{sec:classifier}):
        Immediate practical value as a fast filter in the existing WFST
        pipeline.  Requires the least architectural change.
  \item \textbf{Hybrid WFST+neural} (Section~\ref{sec:hybrid}, Variant~A):
        Builds on the existing WFST and new classifier.  Low risk,
        incremental improvement.
  \item \textbf{Autoregressive LM} (Section~\ref{sec:lm}):
        Independent generative model.  Provides a comparison point and
        may discover novel replicator families.
  \item \textbf{RL fine-tuning} (Section~\ref{sec:rl}):
        Warm-start from the LM, optimize for spread/robustness.
        Highest potential payoff but highest training cost.
  \item \textbf{Masked model} (Section~\ref{sec:masked}):
        Useful for analysis and targeted design.  Lower priority for
        generation but valuable for understanding.
  \item \textbf{SSM} (Section~\ref{sec:ssm}):
        Research exploration---test whether a learned SSM can recover
        and extend the WFST's structure.
  \item \textbf{Seq2seq} (Section~\ref{sec:seq2seq}):
        Requires the richest training data.  Deferred until a large
        library of viable replicators exists.
\end{enumerate}


%% ============================================================
\section{Discussion}
\label{sec:discussion}

The fundamental tension across all approaches is \emph{data efficiency
vs.\ model expressivity}.  The WFST baseline encodes strong domain
knowledge (6502 instruction structure, memory map layout, loop mechanics)
directly into its topology, enabling learning from ${\sim}10^3$ oracle
calls.  Neural approaches trade this inductive bias for representational
flexibility, but require proportionally more data.

The most promising path forward is \emph{progressive neural augmentation}
of the existing WFST pipeline:
\begin{enumerate}
  \item Use the WFST for what it does exactly (structural parsing, proposal
        distribution).
  \item Add a neural surrogate to accelerate filtering.
  \item Optionally replace the WFST's fixed weights with learned,
        context-dependent weights.
  \item Use RL to optimize for objectives beyond viability.
\end{enumerate}

A key open question is whether the viable replicator distribution has
structure beyond what the WFST captures.  If replicator families are
essentially finite and well-separated, the WFST suffices and neural
models add little.  If there is a rich landscape of marginal replicators
with complex interdependencies, neural models will prove essential.

Empirically resolving this question requires growing the corpus of
known replicator families---which itself benefits from the faster
search enabled by the surrogate classifier.  The approaches are thus
mutually reinforcing.


%% ============================================================
\bibliographystyle{plain}
\begin{thebibliography}{99}

\bibitem{chen2021codex}
M.~Chen, J.~Tworek, H.~Jun, et~al.
\newblock Evaluating large language models trained on code.
\newblock \emph{arXiv preprint arXiv:2107.03374}, 2021.

\bibitem{pei2020trex}
K.~Pei, Z.~Xuan, J.~Yang, S.~Jana, and B.~Ray.
\newblock {TREX}: Learning execution semantics from micro-traces for binary
  similarity.
\newblock \emph{arXiv preprint arXiv:2012.08680}, 2020.

\bibitem{lin2017focal}
T.-Y.~Lin, P.~Goyal, R.~Girshick, K.~He, and P.~Doll{\'a}r.
\newblock Focal loss for dense object detection.
\newblock In \emph{Proc.\ ICCV}, pp.~2980--2988, 2017.

\bibitem{snoek2012practical}
J.~Snoek, H.~Larochelle, and R.~P.~Adams.
\newblock Practical {B}ayesian optimization of machine learning algorithms.
\newblock In \emph{Proc.\ NeurIPS}, pp.~2951--2959, 2012.

\bibitem{bieber2020learning}
D.~Bieber, C.~Sutton, H.~Larochelle, and D.~Tarlow.
\newblock Learning to execute programs with instruction pointer attention graph
  neural networks.
\newblock In \emph{Proc.\ NeurIPS}, 2020.

\bibitem{bunel2018leveraging}
R.~Bunel, M.~Hausknecht, J.~Devlin, R.~Singh, and P.~Kohli.
\newblock Leveraging grammar and reinforcement learning for neural program
  synthesis.
\newblock In \emph{Proc.\ ICLR}, 2018.

\bibitem{le2022coderl}
H.~Le, Y.~Wang, A.~D.~Gotmare, S.~Savarese, and S.~C.~H.~Hoi.
\newblock {CodeRL}: Mastering code generation through pretrained models and deep
  reinforcement learning.
\newblock In \emph{Proc.\ NeurIPS}, 2022.

\bibitem{williams1992simple}
R.~J.~Williams.
\newblock Simple statistical gradient-following algorithms for connectionist
  reinforcement learning.
\newblock \emph{Machine Learning}, 8(3):229--256, 1992.

\bibitem{ranzato2016sequence}
M.~Ranzato, S.~Chopra, M.~Auli, and W.~Zaremba.
\newblock Sequence level training with recurrent neural networks.
\newblock In \emph{Proc.\ ICLR}, 2016.

\bibitem{schulman2017ppo}
J.~Schulman, F.~Wolski, P.~Dhariwal, A.~Radford, and O.~Klimov.
\newblock Proximal policy optimization algorithms.
\newblock \emph{arXiv preprint arXiv:1707.06347}, 2017.

\bibitem{tufano2019empirical}
M.~Tufano, C.~Watson, G.~Bavota, M.~{Di Penta}, M.~White, and D.~Poshyvanyk.
\newblock An empirical study on learning bug-fixing patches in the wild via
  neural machine translation.
\newblock \emph{ACM Trans.\ Softw.\ Eng.\ Methodol.}, 28(4):1--29, 2019.

\bibitem{chen2020program}
Y.~Chen, C.~Qi, and G.~Gao.
\newblock Program synthesis with learned code idioms.
\newblock \emph{arXiv preprint}, 2020.

\bibitem{devlin2019bert}
J.~Devlin, M.-W.~Chang, K.~Lee, and K.~Toutanova.
\newblock {BERT}: Pre-training of deep bidirectional transformers for language
  understanding.
\newblock In \emph{Proc.\ NAACL-HLT}, pp.~4171--4186, 2019.

\bibitem{feng2020codebert}
Z.~Feng, D.~Guo, D.~Tang, et~al.
\newblock {CodeBERT}: A pre-trained model for programming and natural languages.
\newblock In \emph{Proc.\ EMNLP Findings}, pp.~1536--1547, 2020.

\bibitem{ghazvininejad2019mask}
M.~Ghazvininejad, O.~Levy, Y.~Liu, and L.~Zettlemoyer.
\newblock Mask-predict: Parallel decoding of conditional masked language models.
\newblock In \emph{Proc.\ EMNLP-IJCNLP}, pp.~6112--6121, 2019.

\bibitem{gu2024mamba}
A.~Gu and T.~Dao.
\newblock Mamba: Linear-time sequence modeling with selective state spaces.
\newblock \emph{arXiv preprint arXiv:2312.00752}, 2023.

\bibitem{mohri2002weighted}
M.~Mohri, F.~Pereira, and M.~Riley.
\newblock Weighted finite-state transducers in speech recognition.
\newblock \emph{Computer Speech \& Language}, 16(1):69--88, 2002.

\bibitem{schwartz2018sopa}
R.~Schwartz, S.~Thomson, and N.~A.~Smith.
\newblock {SoPa}: Bridging {CNN}s, {RNN}s, and weighted finite-state machines.
\newblock In \emph{Proc.\ ACL}, pp.~295--305, 2018.

\end{thebibliography}

\end{document}
